\begin{abstract}

\emph{Self-distillation} has emerged as a promising recipe for self-improvement in language models \citep{zhao2026selfdistilledreasoneronpolicyselfdistillation,shenfeld2026self,hubotter2026reinforcement}.
In this setting, a model can be used as its \textit{own} teacher when augmented with privileged information (e.g. a solution to a math problem). The approach seems especially appealing  for thinking models, which can leverage test-time reasoning to fully absorb the privileged information. Surprisingly,
we show that privileged self-distillation degrades thinking models with respect to long reasoning traces: across five Qwen3 and OLMo thinking models evaluated on AIME24, AIME25, and HMMT25, privileged-context distillation causes a relative drop of up to 17\% in avg@16 accuracy.
The degradation scales with the amount of privileged context withheld from the student and is most pronounced at long rollout budgets, where thinking models otherwise obtain their largest gains.
This failure mode is not specific to self-distillation: on-policy distillation (OPD) improves thinking models, but privileged on-policy distillation reverses these gains.
Our diagnostics suggest that this failure mode is linked to
how privileged teacher context reshapes learning at high-entropy forking positions \citep{bigelow2024forking, zhang2026embarrassingly}, i.e., rollout positions where multiple continuations remain plausible and may lead to different reasoning paths.
Privileged context lowers fork rates in thinking-model rollouts but not in instruction model rollouts.
This leads to an interesting dichotomy wherein privileged context can help instruction-tuned models but hurts more performant thinking models that depend heavily on exploration and rollout quality.
This effect is especially visible when the student begins a self-correction branch, where privileged OPD penalizes sampled reconsideration tokens that vanilla OPD supports.
Thinking models trained with a privileged teacher produce fewer verification, backtracking, and hedging markers, even after length normalization.
These findings indicate that applying self-distillation methods to strong thinking models requires further consideration of token-level signal---especially around tokens related to correction and crucial reasoning steps.

\end{abstract}
